\chapter{Objetivos}
A continuación se detallan los objetivos generales y específicos del oscilador armónico cuántico
\section{Objetivo General}

Desarrollar el cálculo cuántico completo del oscilador armónico, obteniendo explícitamente su espectro energético y sus autofunciones, y analizar su relevancia como modelo fundamental de la mecánica cuántica y sus aplicaciones físicas.


\section{Objetivos específicos}

\begin{itemize}

    \item Formular el oscilador armónico clásico mediante la ley de Hooke y el formalismo hamiltoniano, estableciendo la base conceptual para su posterior cuantización.

    \item Construir el Hamiltoniano cuántico del oscilador armónico a partir del principio de cuantización canónica, introduciendo las relaciones de conmutación fundamentales.

    \item Derivar el espectro energético cuantizado del sistema mediante el formalismo de operadores de creación y aniquilación, demostrando la existencia de niveles equiespaciados y de la energía de punto cero.

    \item Obtener las autofunciones explícitas del oscilador armónico en representación de posición, caracterizando sus propiedades matemáticas (ortonormalidad, paridad y estructura nodal) y su relación con los polinomios de Hermite.

    \item Demostrar la equivalencia formal entre el método algebraico y el método analítico de Schrödinger, validando la consistencia interna del formalismo cuántico.

    \item Analizar la relevancia física e histórica del oscilador armónico cuántico, destacando su papel en vibraciones atómicas, fonones en sólidos y su importancia como modelo base en sistemas bosónicos y computación cuántica.

\end{itemize}

